
\UnnumberedChapter{PHỤ LỤC}

% ============================================
% PHỤ LỤC A. KHAI BÁO SỬ DỤNG AI
% ============================================

\section*{Phụ lục A. Khai báo sử dụng công cụ AI}

Trong quá trình thực hiện bài tập lớn học phần
Nhập môn Học máy với đề tài
``Phân loại cảm xúc từ phản hồi của sinh viên'',
nhóm có sử dụng công cụ ChatGPT của OpenAI
để hỗ trợ một số công việc nghiên cứu,
lập trình và trình bày báo cáo.

Các nội dung có sự hỗ trợ của AI bao gồm:

\begin{itemize}

    \item Giải thích yêu cầu của bài tập lớn,
    tìm hiểu quy trình học máy và các thuật
    toán được sử dụng trong đề tài.

    \item Hỗ trợ xây dựng, rà soát và chỉnh sửa
    mã nguồn trong các Jupyter Notebook,
    bao gồm khám phá dữ liệu, tiền xử lý,
    huấn luyện và đánh giá mô hình.

    \item Hỗ trợ phát hiện và xử lý lỗi Python,
    kiểm tra cách truyền dữ liệu giữa các
    notebook và hạn chế rò rỉ dữ liệu.

    \item Hỗ trợ xây dựng các bảng tổng hợp,
    phân tích chỉ số đánh giá và diễn giải
    các trường hợp mô hình dự đoán sai.

    \item Hỗ trợ soạn thảo, hiệu chỉnh nội dung
    và mã nguồn LaTeX cho báo cáo khoa học.

    \item Hỗ trợ thiết kế sơ đồ quy trình học máy
    và trình bày các kết quả thực nghiệm.

\end{itemize}

Các kết quả thực nghiệm được trình bày
trong báo cáo phải được đối chiếu với
đầu ra thực tế của chương trình trên
bộ dữ liệu do giảng viên chỉ định.

AI không được xem là nguồn xác nhận
tính chính xác của nhãn dữ liệu hoặc
chất lượng dự đoán của mô hình.

Nhóm chịu trách nhiệm kiểm tra và
hiệu chỉnh mã nguồn, xác minh số liệu,
đối chiếu tài liệu tham khảo và
giải thích những phương pháp đã sử dụng.

Những kết luận về nguyên nhân dự đoán sai
được xem xét trên cơ sở kết quả thực nghiệm
và các mẫu dữ liệu được kiểm tra,
không mặc nhiên coi mọi nhận xét do
AI đề xuất là kết luận đã được chứng minh.

Nhóm cam kết khai báo trung thực việc
sử dụng AI và chịu trách nhiệm đối với
nội dung, mã nguồn và kết quả
được nộp trong bài tập lớn.


% ============================================
% PHỤ LỤC B. HƯỚNG DẪN CHẠY CHƯƠNG TRÌNH
% ============================================

\section*{Phụ lục B. Hướng dẫn chạy chương trình}

\subsection*{B.1. Chuẩn bị môi trường}

Chương trình được phát triển bằng Python
và thực thi thông qua Jupyter Notebook
trong Visual Studio Code.

Để cài đặt các thư viện cần thiết,
người sử dụng mở terminal tại thư mục
gốc của dự án và thực hiện:

\begin{verbatim}
python -m pip install -r requirements.txt
\end{verbatim}

Sau đó, lựa chọn môi trường Python
phù hợp làm Jupyter Kernel trong
Visual Studio Code.

\subsection*{B.2. Chuẩn bị dữ liệu}

Tệp dữ liệu đầu vào có tên:

\begin{verbatim}
student_feedback_dataset.csv
\end{verbatim}

Tệp được đặt tại đường dẫn:

\begin{verbatim}
data/student_feedback_dataset.csv
\end{verbatim}

Chương trình cũng cần có module
\texttt{src/text\_utils.py}
để thực hiện tiền xử lý văn bản.

\subsection*{B.3. Thứ tự chạy notebook}

Các notebook được thực thi theo thứ tự:

\begin{enumerate}

    \item \texttt{01\_EDA.ipynb}:
    Khám phá và phân tích dữ liệu.

    \item \texttt{02\_Preprocessing.ipynb}:
    Làm sạch văn bản, chia dữ liệu
    và tạo đặc trưng BoW/TF-IDF.

    \item \texttt{03\_Training.ipynb}:
    Huấn luyện, kiểm định chéo,
    lựa chọn mô hình và dự đoán.

    \item \texttt{04\_Evaluation.ipynb}:
    Đánh giá mô hình,
    xây dựng Confusion Matrix
    và phân tích dự đoán sai.

\end{enumerate}

Mỗi notebook cần được chạy các ô code
theo đúng thứ tự để bảo đảm các biến
và dữ liệu trung gian được khởi tạo.

\subsection*{B.4. Kết quả đầu ra}

Các tệp kết quả được lưu theo
những thư mục chính:

\begin{itemize}

    \item \texttt{results/eda/}:
    Thống kê dữ liệu ban đầu.

    \item \texttt{results/member2\_outputs/}:
    Dữ liệu và đặc trưng sau tiền xử lý.

    \item \texttt{results/training/}:
    Kết quả huấn luyện,
    kiểm định chéo và dự đoán.

    \item \texttt{results/evaluation/}:
    Các bảng đánh giá,
    ma trận nhầm lẫn và phân tích lỗi.

\end{itemize}

Các bước trên cho phép thực hiện lại
quy trình học máy từ dữ liệu đầu vào
đến kết quả đánh giá cuối cùng.


% ============================================
% PHỤ LỤC C. LIÊN KẾT MÃ NGUỒN
% ============================================

\section*{Phụ lục C. Liên kết mã nguồn}

Mã nguồn được nhóm quản lý trên GitHub.

Đường dẫn repository:

\begin{center}
\url{https://github.com/LuuVanNguyen-DNU/Nhom-7_Phan_loai_cam_xuc_tu_phan_hoi_cua_sinh_vien}
\end{center}

Repository cần bao gồm mã nguồn,
các notebook, tệp khai báo thư viện,
README và hướng dẫn chạy chương trình.

Lịch sử commit được sử dụng làm
một trong các minh chứng về quá trình
phối hợp và đóng góp của các thành viên.
